% Appendix B: Full Threat Matrix + Rationality and Scope Blocks
%
% B.1 is the per-design qualitative severity matrix (expanded form of
% Table 5 in Section §3); B.2 is the class rationality and scope
% matrix that populates the seven-field rationality block format of
% AXES.md §7 for each of the six attack classes Z1-Z6. The FC paper
% compresses the rationality blocks to matrix form; the thesis App D
% carries the full free-text form.

\section{Complete Threat Model Matrix}
\label{app:threat_matrix}
\label{app:threat-matrix}

\subsection{Per-Design Severity (TM1--TM11)}
\label{app:threat_matrix:severity}

Table~\ref{tab:pareto-matrix} in Section~\ref{sec:imposs:design_space}
reports worst-case defender loss $\mathcal{L}$ across five dominant
threats (TM1, TM2, TM3, TM4, TM11). Table~\ref{tab:threat-matrix-full}
below lists the remaining threats (TM5--TM10) alongside the dominant
ones for completeness, using a qualitative severity scale
(\textbf{Safe} / \textbf{N/A} / \textbf{Low} / \textbf{Moderate} /
\textbf{High} / \textbf{Severe}) and citing measured vsize where
applicable. Simplicity is excluded from this matrix because
Elements' substrate alters the threat semantics
(Section~\ref{sec:simplicity-perspective}).

\begin{table}[h]
\centering
\caption{Per-design severity for the full 11-element threat
taxonomy. Severity labels: \textbf{N/A} = structurally
inapplicable; \textbf{Low} = requires difficult preconditions;
\textbf{Moderate} = feasible under specific conditions;
\textbf{High} / \textbf{Severe} = significant risk under realistic
assumptions; \textbf{Safe} = correct-by-design; \textbf{Impossible}
= prevented by a cryptographic or structural invariant.}
\label{tab:threat-matrix-full}
\footnotesize
\setlength{\tabcolsep}{3pt}
\begin{tabular}{@{}cp{2.6cm}p{2.6cm}p{2.6cm}p{2.6cm}@{}}
\toprule
TM & \opctv{} & \opccv{} & \opvault{} & \catcsfs{} \\
\midrule
1 & \textbf{High.} 25-descendant chain blocks CPFP (${\approx}14{,}300$~sats).
  & \textbf{N/A} (no anchor).
  & \textbf{N/A} (fee-wallet input).
  & \textbf{N/A} (external fee input). \\[4pt]

2 & \textbf{Moderate.} Hot-key triggers; defender sweeps via cold leaf.
  & \textbf{Severe.} Keyless; $\gamma{=}0.27$.
  & \textbf{Low.} Auth key required; $\gamma{=}0.55$.
  & \textbf{Low.} Cold key required; subsumed by TM11. \\[4pt]

3 & \textbf{High.} Hot+fee $\Rightarrow$ theft; hot alone $\Rightarrow$ DoS.
  & \textbf{High.} Redirect; watchtower races (122\vB).
  & \textbf{Moderate.} Redirect; authorised recovery races (246\vB).
  & \textbf{Low.} Griefing only (dual-verification binding). \\[4pt]

4 & \textbf{N/A} (no partial unvault).
  & \textbf{Severe.} $\sim$6{,}172 splits/block; $r^\dagger \approx 66\satvB$.
  & \textbf{Severe.} $\sim$3{,}427 splits/block; $r^\dagger \approx 59\satvB$.
  & \textbf{N/A} (no partial unvault). \\[4pt]

5 & \textbf{N/A} (no xpub derivation).
  & \textbf{N/A} (no xpub derivation).
  & \textbf{Moderate.} xpub expansion leaks all children.
  & \textbf{N/A} (no xpub derivation). \\[4pt]

6 & \textbf{High.} Hot + fee $\Rightarrow$ theft.
  & \textbf{N/A} (single key; keyless recovery).
  & \textbf{High.} Trigger + auth $\Rightarrow$ persistent DoS.
  & \textbf{High.} Cold key alone $\Rightarrow$ theft (via TM11). \\[4pt]

7 & \textbf{High.} Stuck funds on address reuse.
  & \textbf{Safe.}
  & \textbf{Safe.}
  & \textbf{Safe.} \\[4pt]

8 & \textbf{N/A} (no mode byte).
  & \textbf{Developer footgun.} Undefined mode $\rightarrow$ \texttt{OP\_SUCCESS}.
  & \textbf{N/A} (no mode byte).
  & \textbf{N/A} (no mode byte). \\[4pt]

9 & \textbf{N/A} (no \opcsfs{} path).
  & \textbf{N/A} (no \opcsfs{} path).
  & \textbf{N/A} (no \opcsfs{} path).
  & \textbf{Impossible.} Dual verification rejects redirection. \\[4pt]

10 & \textbf{N/A} (no \opcsfs{} witness).
   & \textbf{N/A} (no \opcsfs{} witness).
   & \textbf{N/A} (no \opcsfs{} witness).
   & \textbf{Impossible.} SHA-256 binding. \\[4pt]

11 & \textbf{N/A} (\opctv{}-committed recovery address).
   & \textbf{N/A} (keyless recovery; no cold key).
   & \textbf{N/A} (Taproot-tweak-constrained recovery).
   & \textbf{High.} Bare \texttt{cold\_pk OP\_CHECKSIG} $\Rightarrow$ theft. \\
\bottomrule
\end{tabular}
\end{table}

\paragraph{Simplicity (cross-substrate reference).}
Under honest-federation assumptions on Elements, TM1 does not
transfer (no public mempool / 25-descendant cap). Our Simplicity
vault is atomic (\axrevault{} = \texttt{no}), so TM4 has no surface;
a partial-withdrawal variant would inherit \classscale{} per the D2
implication (Section~\ref{sec:imposs:class_scale}). TM11 outcome
matches \opctv{} and \opvault{} (output-bound recovery via
\texttt{jet::outputs\_hash()}). TM5, TM8, TM9, TM10 do not apply:
Simplicity has no xpub derivation, no mode byte, no \opcsfs{} path,
and no \opcsfs{} witness. TM7 is safe (per-vault plan). TM6 reduces
to TM11 since the hot path is also output-bound. Full design-space
position and lifecycle measurements:
Appendix~\ref{app:simplicity}.

\subsection{Rationality and Scope Blocks}
\label{app:rationality}

Each attack class \classfee{}--\classmode{} is accompanied by a
seven-field rationality block per the schema of Sections
\ref{sec:imposs:class_fee}--\ref{sec:imposs:class_mode}. For the FC
paper we compress the blocks into Table~\ref{tab:rationality} with
one cell per (class, field) intersection. The thesis Appendix D
carries the uncompressed free-text form.

Every block implicitly cites the Kerckhoffs threat model of
Section~\ref{sec:bg:lifecycle}: the attacker is assumed to know the
vault's design, parameters, and axis values. Where a field is
genuinely N/A for a given class, the cell reads ``N/A --- [reason]''
rather than being omitted.

\begin{table}[t]
\centering
\caption{Rationality and Scope matrix for the six attack classes.
Columns follow the seven-field schema of
Section~\ref{sec:bg:axes}: attacker capability (Cap.), defender
model (Def.), economic preconditions (Econ.), protocol/network
assumptions (Prot.), scope of validity (Scope), rational-attacker
check (Rat.), counter-factual mitigations (CF). Abbreviations: WT =
watchtower; RBF3 = BIP-125 Rule 3; \texttt{dust} = BIP-140/P2TR
$330$-sat or P2WPKH $546$-sat threshold.}
\label{tab:rationality}
\tiny
\setlength{\tabcolsep}{2pt}
\renewcommand{\arraystretch}{1.25}
\begin{tabular}{@{}p{0.9cm}p{1.55cm}p{1.55cm}p{1.55cm}p{1.55cm}p{1.55cm}p{1.55cm}p{1.55cm}@{}}
\toprule
Class & Cap. & Def. & Econ. & Prot. & Scope & Rat. & CF \\
\midrule
\classfee{} (TM1) &
  Hot key; anchor-desc. capital ($\approx 14{,}300$~sats). &
  Online WT; Core default mempool. &
  $r$-invariant; TM1 cost/$V$ $< 3\%$ at any $r$. &
  BIP-125 25-desc. cap; $3\satvB{}$ Core default dustRelay. &
  Applies iff \axfee{}\,=\,\texttt{out-of-band-anchor}; N/A for dynamic fee models. &
  Rational at any $r$ given attacker has hot key; yield $\approx V$. &
  TRUC/BIP-431 collapses desc.\,cap to 1 --- eliminates class; cluster-mempool: scope-limit.
\\
\classgrief{} (TM2) &
  \textbf{Zero keys} for \opccv{}; recovery-address public under Kerckhoffs. &
  Any; attack is pure liveness denial. &
  Linear in $r$; $\gamma(D)$ $r$-invariant. &
  Mempool-observable trigger tx; no RBF dependence. &
  Applies iff \axauth{}\,=\,\texttt{keyless}; reduces to TM6 otherwise. &
  Irrational without external incentive (no financial gain). &
  \axauth{}\,=\,\texttt{recoveryauth-key} (OPV) or key-gated (CTV/\catcsfs) eliminates class.
\\
\classscale{} (TM4) &
  Trigger key; capital for $N$ partial-unvaults. &
  WT with recovery authority; batching flag $b$; minimal-disruption policy maximises attacker leverage. &
  Rate-dep.; threshold $r^\dagger_{\mathrm{TM4}}$ per Eq.~\ref{eq:r-dagger}. &
  Dust thresholds; RBF3; sibling-eviction disallowed; attacker's revault chain is single-chain. &
  Applies iff \axrevault{}\,=\,\texttt{yes}; structural cost tail independent of adversary. &
  Rational only above $r^\dagger$ with external incentive; pure liveness attacker irrational at any $r$. &
  Defender batching: $r^\dagger$ $\times 1.8$--$2.7$; full-sweep-on-first-alert: 1 round only.
\\
\classhot{} (TM3) &
  Hot key; optional fee-pinning capital (if \classfee{} composes). &
  Online WT able to race during CSV; race bias toward attacker under TM1 composition. &
  Race outcome $r$-independent except via TM1. &
  CSV delay; mempool propagation (Ria23 cycling out of scope). &
  Applies iff \axwdbind{}\,$\neq$\,\texttt{signature-dual-bind}. &
  Rational only paired with fee-pinning budget; otherwise defender wins race generically. &
  \axwdbind{}\,=\,\texttt{signature-dual-bind} (CAT+CSFS): Schnorr EUF-CMA forces output-binding --- eliminates class.
\\
\classcold{} (TM11) &
  Cold key (Kerckhoffs). &
  Any; attack instant, no CSV window, no WT intervention. &
  $r$-invariant. &
  N/A --- pure signature check; no mempool dynamics. &
  Applies iff \axrecbind{}\,=\,\texttt{plain-signature}. &
  Rational whenever cold key is compromised. &
  \axrecbind{}\,=\,\texttt{signature-dual-bind} via Poelstra's \catcsfs{} recovery-leaf sketch: eliminates class.
\\
\classmode{} (TM8) &
  Any observer; no keys. &
  Developer discipline around mode-byte selection. &
  N/A --- full vault bypass at constant $110\vB{}$/spend. &
  BIP-341/342 \texttt{OP\_SUCCESSx} semantics: undefined mode terminates script with success. &
  Applies iff \axopsurf{}\,=\,\texttt{multi-mode-with-OP\_SUCCESS}. &
  Rational whenever a misconfigured vault exists on-chain. &
  \axopsurf{}\,=\,\texttt{single-mode}; bitmask-gated mode check; pymatt \texttt{CCV\_FLAG\_*} constants.
\\
\bottomrule
\end{tabular}
\end{table}

\providecommand{\todo}[2][]{\textbf{[TODO: #2]}}
